\subsection{半自反空间}

定义 2. --- 设 E 是一个局部凸空间。若从 E 到 E'' 中的典范映射 \(c_{\mathbf{E}}\) 是双射，则称 E 是半自反的。

这蕴含 E 是 Hausdorff 的，且 \(E'\) 上每个对于强拓扑 \(\beta(E', E)\) 连续的线性型都具有形式 \(x' \mapsto \langle x, x' \rangle\)（其中 \(x \in E\)），即对于弱拓扑 \(\sigma(E', E)\) 连续。

定理 1. --- 局部凸 Hausdorff 空间是半自反的，当且仅当 E 的每个有界子集对于弱化拓扑 \(\sigma(E, E')\) 是相对紧的。若 E 是半自反的，则 E 的强对偶 \(E_b'\) 是桶型的。

第二个断言由（IV，第 15 页的）命题 3 以及 E 对于初始拓扑与对于弱化拓扑的有界子集的一致性（III，第 27 页，推论 3）得出。

说 E 是半自反的，意味着 \(E_{b}^{\prime}\) 上的拓扑与 E 和 \(E^{\prime}\) 之间的对偶相容，换言之，由 Mackey 定理（IV，第 2 页，定理 1），\(E_{b}^{\prime}\) 上的拓扑比 \(\tau(\mathrm{E}^{\prime}, \mathrm{E})\) 粗（且事实上与它相同）；由定义（IV，第 2 页），这就是说 E 的每个闭、凸且有界的子集对于 \(\sigma(\mathrm{E}, \mathrm{E}^{\prime})\) 紧，而这等价于说 E 的每个有界子集对于 \(\sigma(\mathrm{E},\mathrm{E}^{\prime})\) 相对紧，因为 E 的有界子集的闭凸包是有界的（III，第 3 页，命题 1）。

推论. --- 设 E 是一个局部凸半自反空间。E 的每个闭向量子空间 M 都是半自反的；此外，\(E^{\prime}/M^{\circ}\)（视为 M 的对偶）上的强拓扑是 \(E^{\prime}\) 上强拓扑的商拓扑。

设 B 是 M 的一个有界子集。由于 B 在 E 中有界，且 M 上的弱化拓扑 \(\sigma(\mathbf{M},\mathbf{M}^{\prime})\) 由 \(\sigma(\mathbf{E},\mathbf{E}^{\prime})\) 诱导（IV，第 10 页，命题 11），B 在赋以 \(\sigma(\mathbf{M},\mathbf{M}^{\prime})\) 的 M 中的闭包是紧的。于是由定理 1，M 是半自反的。推论的最后一个断言由 IV 第 9 页的命题 10（应用到 E 的所有闭、凸且有界子集组成的集 S）得出。

注. --- 1) 假定 E 是半自反的。E 的每个对于初始拓扑凸、闭且有界的子集对于拓扑 \(\sigma(E, E')\) 是紧的（IV，第 1 页，命题 1）。* 另一方面，无限维 hilbertian 空间 E 的单位球面（以方程 \(\|x\|=1\)）对于初始拓扑是闭且有界的，但对于弱化拓扑不是闭的，即使 E 是半自反的。*

\begin{enumerate}
\def\labelenumi{\arabic{enumi})}
\setcounter{enumi}{1}
\item
  由 IV 第 5 页的注 3，我们可以把定理 1 改述如下：Hausdorff 空间 E 是半自反的，当且仅当它对于其弱化拓扑是拟完备的。若它是半自反的，则它对于其初始拓扑是拟完备的（IV，第 5 页，注 2）。
\item
  在上述推论的假设下，空间 E/M 未必是半自反的（IV，第 63 页，习题 10）。
\end{enumerate}

